\documentclass[12pt,a4paper]{article}
\usepackage[margin=1in]{geometry}
\usepackage{booktabs}
\usepackage{setspace}
\usepackage{amsmath,amssymb}
\usepackage{enumitem}
\usepackage{hyperref}
\setstretch{1.15}

\begin{document}

\begin{center}
{\Large\bfseries Response to Reviewers}\\[2pt]
Manuscript ID: FINEL-D-26-00567\\[2pt]
\textit{Finite Element Analysis of Second-Order Strain-Gradient Elasticity with Ellipsoidal Nonlocal Characteristic Lengths}
\end{center}

\vspace{0.3cm}

We thank the reviewers for their detailed comments. In addition to addressing all reviewer points, we have corrected a mass-matrix regularization error discovered during final verification. The corrected results are described below and reflected in the revised manuscript with blue highlighting.

\textbf{Important correction:} In the original 3D solver \texttt{solve\_3d\_sg\_rotated}, the mass matrix was regularized as $M_{ff}^{reg}(i,i)=M_{ff}(i,i)+10^{-8}|K_{ff}(i,i)|$ for derivative DOFs, based on the assumption that $M_{ff}$ is singular. The consistent mass $M=\int \rho N^T N d\Omega$ is in fact positive-definite (min eigenvalue $1.34\times10^{-13}$, max $12.15$). The regularization multiplied physical mass by $257\times$--$1354\times$ (median $575\times$) on 756 DOFs, causing $f_1\propto \epsilon^{-0.5}$ scaling and violating the bound $f_1>171.61$ Hz (classical frequency on same $4\times2\times2$ mesh). The original orientation sweep gave $f_1=113.3,88.6,69.4,66.4,71.5,79.8,89.1$ Hz with a spurious minimum at $45^\circ$. After correction $M_{scaled}=S_{scale}M_{ff}S_{scale}$ with $eigs(K_{scaled},M_{scaled},6,'smallestabs')$ and dense fallback, the corrected sweep gives $f_1=186.4,185.6,183.5,180.5,177.3,174.6,173.4$ Hz, $f_{iso}=177.0$ Hz ($f_{iso}=178.72$ Hz for $l_{ref}=0.05$ m reference in Fig.~11), $u/u_0=1.00\to1.039$, monotonic decrease, no minimum at $45^\circ$. The isotropic reference $u_{ref}=4.7876\times10^{-6}$ m matches $EA/L$ estimate. Static results are unchanged ($u$ ratio $1.00\to1.04$). All 3D figures have been regenerated with the corrected solver; font Times New Roman and tiledlayout are preserved. The theoretical conclusion about orientation-dependent tuning remains, but the non-monotonic minimum at $45^\circ$ is removed as an artifact of the regularization.

% ============================================================
% REVIEWER 1
% ============================================================
\section*{Reviewer 1}

\begin{enumerate}[leftmargin=14pt, itemsep=10pt]

\item \textbf{Comment:} The core scientific contributions and new concepts must be clearly highlighted in both the Abstract and the Introduction sections.

\textbf{Response:} Abstract revised to four-part narrative and five contributions listed: (1) geometric derivation of $L_{mn}$ from second moments of ellipsoid, (2) positive-definite anisotropic energy, (3) unified $C^1$ FEM, (4) orientation-dependent monotonic tuning governed by $L_{11}(\theta)$ and $L_{12}$ coupling, (5) bivariate AR-$\theta$ design map and DSI.

\item \textbf{Comment:} The Abstract should be restructured to provide a concise narrative covering the primary research question, methodology, key findings, and practical engineering significance.

\textbf{Response:} Revised Abstract now follows this structure, with corrected values $f_1=186.4\to173.1$ Hz, $f_{iso}=178.7$ Hz, $u_{ref}=4.79$ $\mu$m, $f_1/f_{iso}=1.054\to0.979$.

\item \textbf{Comment:} Clarify the treatment of nonlocal domain truncation near physical boundaries and its impact on higher-order boundary conditions.

\textbf{Response:} Added explicit statement in Sec.~2.1 that derivation assumes $\mathcal V$ fully inside body, bulk approximation, position-dependent $L_{mn}(\mathbf x)$ near boundaries neglected as second-order for $l_{max}\ll dist(\mathbf x,\partial\Omega)$. Future work noted.

\item \textbf{Comment:} Justify retaining the 3D factor (1/10) in 1D and 2D formulations rather than re-integrating over 1D segments or 2D ellipses.

\textbf{Response:} Appendix A added with derivation: 1D $l^2/3\to1/6$, 2D $l^2/4\to1/8$, 3D $l^2/5\to1/10$. Retain $1/10$ because physical neighbourhood remains 3D ellipsoid, 1D/2D are reduced kinematics.

\item \textbf{Comment:} Demonstrate how a coarse 3D mesh (4 x 2 x 2) accurately captures steep strain-gradient boundary layers at clamped ends without inducing artificial stiffening.

\textbf{Response:} For $4\times2\times2$ mesh at $l/L_x=0.05$, $u_{avg}=4.7876$ $\mu$m, reaction error $1.06\times10^{-12}\%$, $\kappa=3.27\times10^5$. Tricubic Hermite carries derivative DOFs, clamped-gradient enforced essentially, $3\times3\times3$ Gauss exact, diagonal scaling keeps $\kappa<10^6$.

\item \textbf{Comment:} Explain the process for extending this tensor-product $C^1$ Hermite element framework to irregular domain geometries or isoparametric meshes.

\textbf{Response:} Discussed isoparametric $C^1$ Hermite with Jacobian/Hessian, and alternatives Argyris, Bell, IGA, VEM.

\item \textbf{Comment:} Provide mode shape visualizations to illustrate the physical energy redistribution driving the frequency minimum at $\theta = 45^0$.

\textbf{Response:} Original minimum at $45^\circ$ was artifact of $M_{ff}^{reg}$ regularization as explained above. With corrected mass, $L_{11}=0.04\to0.0004$, $L_{12}=-0.0198$ max at $45^\circ$, but $f_1$ decreases monotonically $186.4\to173.1$ Hz, no minimum. Energy redistribution still occurs via $L_{12}B_{,x}^TCB_{,y}$ but does not produce softening below classical bound. Mode shapes remain smooth.

\item \textbf{Comment:} Outline a practical experimental protocol to calibrate the characteristic length parameters ($l_1, l_2, l_3$) for real heterogeneous materials.

\textbf{Response:} Added Sec.~6.1 six-step protocol: bending/torsion size-effect, directional tests, ultrasonic dispersion, DIC boundary layer, RVE homogenization, validation on AM lattices.

\item \textbf{Comment:} Justify the omission of strain-gradient kinetic energy (micro-inertia) from the dynamic formulation.

\textbf{Response:} Without micro-inertia $\omega^2=c^2k^2(1+l_s^2k^2/10)$. Including $T_g$ gives $\omega^2=c^2k^2(1+l_s^2k^2/10)/(1+l_k^2k^2)$. For $lk<0.3$ for first modes, effect $<5\%$, omitted for low-frequency focus.

\item \textbf{Comment:} Detail how higher-order line forces at sharp corners and edges are resolved.

\textbf{Response:} Full Mindlin-Toupin has edge forces $e_i=[\![q_{ijm}\nu_m\mu_j]\!]$. Simplified Aifantis-type uses reduced $q_{ijm}\nu_m\nu_j=0$, Hermite DOFs shared at corners enforce variationally, error $<10^{-11}$.

\item \textbf{Comment:} Define a physical phenomenon not extensively studied.

\textbf{Response:} Orientation-dependent monotonic tuning governed by $L_{11}(\theta)$ and $L_{12}$ coupling, with three design regions in AR-$\theta$ space, is highlighted as new contribution after correction.

\end{enumerate}

% REVIEWER 2
\section*{Reviewer 2}

\begin{enumerate}[leftmargin=14pt, itemsep=10pt]

\item \textbf{Comment:} Nonlocal parameter as optimization variable?

\textbf{Response:} $l_i$ are material constants, AR and $\theta$ are processing-induced design manifestations, not actively tuned with fields. Calibration protocol in Sec.~6.1.

\item \textbf{Comment:} $l_{max}\ll L_\varepsilon$ vs $l/L$ up to 0.20 and $l_x=1.0$ m?

\textbf{Response:} Larger values are parametric sensitivity study, not strict weakly nonlocal. Main conclusions drawn for $l/L\le0.10$, stated explicitly.

\item \textbf{Comment:} Taylor expansion eq. (9) gives $\tilde\epsilon=\epsilon+(1/10)L\epsilon$, energy coefficient 1/20 vs 1/10?

\textbf{Response:} $1/20$ in energy gives $1/10$ in double stress after variation due to symmetry, derivation added in Sec.~2.3.

\item \textbf{Comment:} Non-classical boundary conditions?

\textbf{Response:} Essential $u_{i,\nu}=\bar g_i$ via Hermite DOFs, natural $q_{ijm}\nu_m\nu_j=0$ variationally.

\item \textbf{Comment:} Double traction orientation dependence at $\theta=45^\circ$?

\textbf{Response:} $r_{ij}=1/10[L_{11}\sigma^{cl}_{ij,1}+L_{12}\sigma^{cl}_{ij,2}]$ on $x=0$ face, $L_{12}$ max at $45^\circ$ couples transverse gradient, physical consequence of rotated anisotropy, retained in reduced BC as normal projection.

\item \textbf{Comment:} Boundary truncation for thin structures?

\textbf{Response:} Position-dependent $L_{mn}(\mathbf x)$ needed when $l\sim h$, full integral model preferable, noted as future work.

\item \textbf{Comment:} Introduction brief, extend with six references.

\textbf{Response:} Added paragraph reviewing Phung-Van et al. AMM 128, AMM 127, CS 335, EC 39, TWS 198, and related works, with corrected citations.

\item \textbf{Comment:} No comparison against full integral nonlocal, homogenized lattice, alternative gradient.

\textbf{Response:} Added discussion: differential model is $O(l^4)$ approximation to integral Gaussian kernel, $L_{mn}$ corresponds to RVE second moments, comparison with Mindlin Form II/III.

\item \textbf{Comment:} Steel lengths unrealistic, how calibrate $l_i$?

\textbf{Response:} Steel $l\sim\mu$m, $l/L=0.20$ exaggerated for illustration; fibre composite $l_x$ mm-cm, $l_y$ $\mu$m, AR 10-20 realistic; AM lattice $l\sim$ cell size. Calibration via bending, RVE.

\end{enumerate}

% REVIEWER 3
\section*{Reviewer 3}

\begin{enumerate}[leftmargin=14pt, itemsep=10pt]

\item \textbf{Comment:} Physics of anisotropy unclear, local isotropic but ellipsoidal domain.

\textbf{Response:} Two mechanisms: $C_{ijkl}$ anisotropy vs $L_{mn}$ anisotropy. Present work isolates second mechanism relevant for fibre composites, drawn metals, AM lattices where matrix isotropic but interaction range direction-dependent. $L_{mn}$ analogous to fabric tensor, established in Lazar-Po 2015, Polizzotto 2018, Auffray 2019.

\item \textbf{Comment:} Connection Eq.(7) and Eq.(12).

\textbf{Response:} Eq.(7) kinematic Taylor result, Eq.(12) constitutive Aifantis-type choice ensuring positive-definiteness and $1/10$ in stress after variation, clarified in Sec.~2.3.

\item \textbf{Comment:} Numerical examples do not show practical value beyond isotropic.

\textbf{Response:} Axial vs transverse difference $2.3\%$ for $(0.10,0.02)$ vs $(0.02,0.10)$, $5.3\%$ for $(0.20,0.02)$ vs $(0.02,0.20)$; axial-dominant $5.8\%$ vs transverse $0.5\%$; orientation monotonic $1.00\to1.04$ static and $186.4\to173.1$ Hz dynamic; AR-$\theta$ map reveals stiffening-optimal and insensitive regions, DSI quantifies.

\item \textbf{Comment:} Parameter range vs weakly nonlocal.

\textbf{Response:} Larger values sensitivity study, main conclusions $l/L\le0.10$, stated explicitly.

\item \textbf{Comment:} Wave-dispersion analysis contribution?

\textbf{Response:} Stability check $1+l^2k^2/10>0$, short-wavelength stiffening, validity $lk\ll1$, retained as analytical check.

\item \textbf{Comment:} FEM formulations repetitive, coefficient 1/10.

\textbf{Response:} Restructured Sec.~4 starting with general 3D weak form, 1D/2D as reduced specialisations, coefficient justification in Appendix A.

\item \textbf{Comment:} 3D verification coarse mesh.

\textbf{Response:} $4\times2\times2$ 45 nodes 1080 DOF, $u_{avg}=4.7876$ $\mu$m, $f_{iso}=178.72$ Hz $>171.61$ Hz classical, $\kappa=3.27\times10^5$, reaction error $<10^{-11}$, mesh sensitivity expected $<0.3\%$.

\item \textbf{Comment:} Language editing.

\textbf{Response:} Thorough edit done, terminology unified to characteristic length $l_i$, captions rewritten, no ``to be placed'' placeholders remain.

\end{enumerate}

We believe the revised manuscript with corrected mass formulation and monotonic orientation results is now consistent, numerically verified, and free of placeholder text.

\end{document}
